\subsection{Heurística}

\texttt{manhattanHeuristic} ya viene provista en \texttt{searchAgents.py} (no fue necesario
implementarla para esta actividad):

\begin{verbatim}
def manhattanHeuristic(position, problem, info={}):
    xy1 = position
    xy2 = problem.goal
    return abs(xy1[0] - xy2[0]) + abs(xy1[1] - xy2[1])
\end{verbatim}

Calcula $h_M(n) = |x_n - x_g| + |y_n - y_g|$: la suma de las diferencias absolutas en cada eje,
que es exactamente el número mínimo de movimientos ortogonales (norte/sur/este/oeste) necesarios
para llegar de $n$ a la meta \emph{si no hubiera paredes de por medio}.

\subsection{Procedimiento y resultados}

Se ejecutó A* con Manhattan sobre \texttt{mediumClassic} y se comparó contra UCS
(ver \texttt{experimentos/actividad5\_manhattan.py}). En vez de una interfaz gráfica (no disponible
en el entorno de ejecución usado), se registró el conjunto de celdas efectivamente expandidas por
cada algoritmo (\texttt{problem.\_visitedlist}, la misma estructura que \texttt{graphicsDisplay.py}
usaría para pintarlas en pantalla).

\begin{table}[H]
\centering
\caption{UCS vs. A* con distancia Manhattan}
\label{tab:act5-manhattan}
\begin{tabular}{lrrr}
\toprule
\textbf{Algoritmo} & \textbf{Costo} & \textbf{Expandidos} & \textbf{Tiempo (s)} \\
\midrule
UCS & 12 & 69 & 0.000251 \\
A* + Manhattan & 12 & 15 & 0.000068 \\
\bottomrule
\end{tabular}
\end{table}

Ambos encuentran el mismo costo óptimo (12), pero A*+Manhattan expande 4.60 veces menos nodos que
UCS (69 → 15): de las 70 celdas que UCS llegó a expandir, 54 nunca fueron tocadas por A*+Manhattan.
Esto es evidencia directa (no solo teórica) de que la heurística está orientando la búsqueda hacia
el objetivo en lugar de explorar por igual en todas las direcciones.

\subsection{Pregunta de análisis}

\textbf{¿Por qué la distancia Manhattan puede orientar correctamente a Pac-Man aunque la
heurística no considere directamente las paredes del laberinto?}

Manhattan es admisible precisamente porque \emph{ignora} las paredes: al no considerarlas, nunca
puede sobreestimar el costo real, porque el camino real (que sí tiene que rodear paredes) nunca
puede ser más corto que la línea de movimientos ortogonales sin obstáculos que calcula $h_M$. Es
decir, $h_M(n)$ es siempre una \emph{cota inferior} del costo real $h^*(n)$: en el peor caso (un
laberinto muy enredado) subestima mucho y aporta poca guía, pero nunca lleva a A* a expandir un
nodo equivocado creyendo que está más cerca de lo que en realidad está.

Además, aunque no vea las paredes explícitamente, Manhattan sigue \emph{ordenando} correctamente
los estados por qué tan lejos están de la meta en la métrica correcta para este problema: como
Pac-Man solo se mueve en las 4 direcciones cardinales (nunca en diagonal), la distancia Manhattan
es exactamente la noción de ``cercanía'' que corresponde a este espacio de acciones. Eso es lo que
le permite priorizar, entre varios sucesores con el mismo $g(n)$, el que está en la dirección
correcta hacia la meta, incluso cuando localmente esa dirección choque con una pared unos pasos
más adelante (en ese caso, A* simplemente termina expandiendo también el nodo que rodea la pared,
porque su $g(n)$ real sigue contando el costo verdadero recorrido).
